\subsection{Variazioni dipendenti dal tempo}
In questo caso si suppone che gli stati rimangano gli stessi con lo scorrere del tempo, e la variazione opera cambiando le densità di probabilità dei livelli.

\subsection*{Rappresentazione di interazione}
Si parte dallo stato generico $\Ket{\psi(t)}$ che segue l'andamento determinato dall'Hamiltoniana totale $\H(t) = \H_0 + V(t)$, con $\H_0$ l'hamiltoniana non perturbata e $V(t)$ la perturbazione dipendente dal tempo. Chiaramente lo stato evolve nel modo determinato dall'operatore di evoluzione temporale:
\[
    \Ket{\psi(t)} = U(t,t_0) \Ket{\psi(t_0)} = e^{-\frac{i}{\hslash} \H_0 (t-t_0)}\Ket{\psi_I(t)} 
\]
dove si è definito lo stato nell'interazione come:
\[    \Ket{\psi_I(t)} = e^{\frac{i}{\hslash} \H_0 (t-t_0)} \Ket{\psi(t)} \]
Inserendo questa definizione nell'equazione di Schrödinger si ottiene, sapendo che $\Ket {\psi_I(t_0)} = \Ket {\psi(0)}$:
\[
    i \hslash \tderiv{} \left( e^{-\frac{i}{\hslash} \H_0 (t-t_0)}\Ket{\psi_I(t)} \right) = \H(t) e^{-\frac{i}{\hslash} \H_0 (t-t_0)}\Ket{\psi_I(t)} \]\[
    i \hslash \tderiv{} \Ket{\psi_I(t)} = V_I(t) \Ket{\psi_I(t)}
\]
con $V_I(t) = e^{\frac{i}{\hslash} \H_0 (t-t_0)} V(t) e^{-\frac{i}{\hslash} \H_0 (t-t_0)}$. Questa fattorizzazione è utile in quanto risulta molto più semplice dell'evoluzione temporale fatta direttamente con l'Hamiltoniana perturbata.

\subsubsection*{Sviluppo di Dyson}
A questo punto si può procedere come segue:
\[
    \left[ \Ket {\psi_I(t)} - \Ket {\psi_I(0)} \right] = \frac{1}{i \hslash} \int_0^t V_I(t') \Ket {\psi_I(t')} dt' \]\[
    \Ket {\psi_I(t)} = \Ket {\psi_I(0)} + \frac{1}{i \hslash} \int_0^t V_I(t') \Ket {\psi_I(t')} dt'
\]
in questo modo si rappresenta l'evoluzione dello stato d'interazione in termini dello stato iniziale più un'integrale, espandibile nel seguente modo, in maniera ricorsiva:
\[
    \Ket {\psi_I(t)} = \Ket {\psi_I(0)} + \frac{1}{i \hslash} \int_0^t V_I(t') \Ket {\psi_I(0)} dt' + \left( \frac{1}{i \hslash} \right)^2 \int_0^t dt' \int_0^{t'} dt'' V_I(t') V_I(t'') \Ket {\psi_I(0)} + \ldots
\]
Questa espansione prende il nome di \textbf{sviluppo di Dyson}\index{Sviluppo di Dyson}.
Per linearità si può comodamente scrivere:
\[
    \Ket {\psi_I(t)} = \left[ \hat {Id} + \frac{1}{i \hslash} \int_0^t V_I(t') dt' + \left( \frac{1}{i \hslash} \right)^2 \int_0^t dt' \int_0^{t'} dt'' V_I(t') V_I(t'') + \ldots \right] \Ket {\psi_I(0)}
\]
definendo quindi un operatore di evoluzione nell'interazione tramite il T-prodotto esponenziale (\textbf{operatore di Dyson}\index{Operatore di Dyson}), che in simboli si scrive
\[
    \U_I(t) = T - \exp \{ -\frac{i}{\hslash} \int_0^t V_I(t') dt' \}
\]

È importante notare che non si tratta di uno sviluppo di Taylor, in quanto i termini non sono semplicemente potenze successive di un operatore, ma contengono integrali temporali ordinati. Tuttavia se $[V_I(t'), V_I(t'')] = 0$ per ogni $t', t''$, comparirebbe un termine fattoriale al denominatore, e allora l'operatore di Dyson si riduce a una semplice esponenziale.

\subsection*{Evoluzione temporale degli operatori}
Un operatore evolve nel tempo secondo la solita legge:
\[
    \hat A_I(t) = U^\dagger(t,t_0) \hat A_0 U(t,t_0) = e^{\frac{i}{\hslash} \H (t-t_0)} \hat A_0 e^{-\frac{i}{\hslash} \H (t-t_0)}
\]
E la sua media temporale è
\[
    \avg[t]{\hat A} = \Bra{\psi(t)} \hat A_0 \ket{\psi(t)} = \Bra{\psi_I(t)} \hat A_I(t) \ket{\psi_I(t)}
\]
Si calcoli per esercizio l'equazione di Heisenberg per lo stato $A_I(t)$ (soluzione: \( \frac{i}{\hslash} \tderiv{} \hat A_I(t) = [\hat A_I(t), \H_0] \), ovvero l'equazione di Heisenberg senza termine di perturbazione).



\subsubsection*{Evoluzione delle componenti di uno stato}
Si consideri uno stato generico $\Ket{\psi(t)}$. La sua rappresentazione d'interazione $\Ket{\psi_I(t)}$ può essere espressa nella base degli autostati di $\H_0$, espreimendo l'evoluzione temporale delle componenti:
\[
    \Ket{\psi_I(t)} = \sum_n c_n(t) \Ket{u_n^{(0)}}\ ,\ \ C_n(t) = \Bra{u_n^{(0)}} \Ket{\psi_I(t)}
\]
allora l'equazione di Schrödinger nell'interazione diventa:
\[
    i \hslash \tderiv{} \sum_n C_n(t) \Ket{u_n^{(0)}} = V_I(t) \sum_n C_n(t) \Ket{u_n^{(0)}}
\]
Moltiplicando a sinistra per $\Bra{u_m^{(0)}}$ si ottiene:
\[
    i \hslash \tderiv{} C_m(t) = \sum_n V_{m,n}(t) C_n(t)
\]
con $V_{m,n}(t) = \Bra{u_m^{(0)}}| V_I(t) |u_n^{(0)}\ket$ (secondo la rappresentazione per cui i $V_{m,n}(t)$ sono gli elementi di matrice della perturbazione nell'interazione).
E quindi:
\[
    i\hslash \tderiv {C_n(t)} = \sum_m V_{n,m}(t) C_m(t) e^{i \omega_{n,m} t}
\]
che è un sistema di equazioni differenziali lineari!


Riscriviamo ora i $C_n(t)$ in termini dello stato imperturbato:
\[
    C_n(t) = \braket{n}{\psi_I(t)} = \braket{n}{e^{\frac{i}{\hslash}\H_0 t}\U(t)\psi(t)}
\]
Leoleo

\subsection*{Oscillazioni di Rabi}
Le oscillazioni di Rabi descrivono la dinamica coerente di popolazione in un sistema a due livelli \(\{\Ket{g},\Ket{e}\}\) soggetto a un campo di guida armonico. Con detuning \(\Delta=\omega-\omega_0\) (differenza tra frequenza del campo \(\omega\) e la transizione libera \(\omega_0\)) e ampiezza di accoppiamento caratterizzata dalla Rabi-frequency \(\Omega\) (proporzionale all'intensità del campo e alla matrice di dipolo), nell'approssimazione dell'onda rotante (RWA) l'Hamiltoniana efficace è
\[
    H_{\rm eff} = -\frac{\hbar}{2}\big(\Delta\sigma_z + \Omega\sigma_x\big).
\]
Le probabilità di trovare il sistema nello stato eccitato partendo da \(\Ket{g}\) sono
\[
    P_e(t)=\frac{\Omega^2}{\Omega_{\rm R}^2}\sin^2\!\Big(\frac{\Omega_{\rm R} t}{2}\Big),
    \qquad
    \Omega_{\rm R}\equiv\sqrt{\Delta^2+\Omega^2}
\]
dove \(\Omega_{\rm R}\) è la Rabi-frequency effettiva (inclusa la detuning). In particolare, in risonanza \((\Delta=0)\) si ottiene inversione completa di popolazione e
\[
    P_e(t)=\sin^2\!\Big(\frac{\Omega t}{2}\Big),
\]
permettendo operazioni controllate come i pulse \(\pi\) e \(\pi/2\). Effetti dissipativi (decoerenza, damping) smorzano queste oscillazioni introducendo decadimento esponenziale delle ampiezze.

Si parte da un Hamiltoniana $H_0$ che ammette base formata da due stati:
\[
    H_0 = E_0 \ketbra{0}{0} + E_1\ketbra{1}{1} 
\]
e un potenziale perturbativo
\[
    V(t) = \gamma e^{i\omega t}\ketbra{0}{1} + \gamma^* e^{-i\omega t}\ketbra{1}{0}
\]
che è rappresentabile come
\[
V(t)=\begin{pmatrix}
0 & \gamma e^{i\omega t}\\
\gamma^* e^{-i\omega t} & 0
\end{pmatrix}
\]
quindi, dato uno stato generico $\Ket \psi = C_0\Ket 0 + C_1 \Ket 1$, l'equazione di Schrödinger è:
\[
    i\hslash \tderiv{}
    \begin{pmatrix}
    C_0\\ C_1
    \end{pmatrix}
    = \begin{pmatrix}
    0 & \gamma e^{i(\omega + \omega_{01}) t}\\
    \gamma^* e^{-i(\omega + \omega_{01}) t} & 0
    \end{pmatrix}
    \begin{pmatrix}
    C_0\\ C_1
    \end{pmatrix}
\]
Partiamo ponendo $\Ket {\psi_0} = \Ket 1$. Il sistema di equazioni differenziali è risolvibile, e risulta:
\[
\begin{cases}
    i\hslash \dot C_0 =  \gamma e^{i(\omega + \omega_{01}) t} C_1 \\
    i\hslash \dot C_1 =  \gamma^* e^{-i(\omega + \omega_{01}) t} C_0 \\
\end{cases}
\]
al chè (si completico i calcoli per esercizio):
\[
    C_0(t) = i\left[e^{i(\omega - \omega_{21})t} \left(A\sin(\Omega t) + B\sin (\Omega t)\right)\right]
\]
E risulta che $B = 1$ e $A = \frac{-1(\omega - \omega_{21})}{2\Omega}$. Si calcola la probailità che uno stato $\Ket {\psi(t)} = \Ket 0$ rimanga nello stato 0:
\[
    |C_0(t)|^2 = \cos^2(\Omega t) ^ \frac{(\omega - \omega_{21})^2}{4\Omega^2}\sin^2(\Omega t)
\]
e la probabilità di transizione è quindi 

\[
    |C_1(t)|^1 = 1 - |C_=(t)|^2 = \sin^2{\Omega t}\left(\frac{4\Omega - (\omega - \omega_{21})^2}{4\Omega^2}\right)
\]
ovvero una probailità che oscilla periodicamente nel tempo. In condizione di risonanza ($\omega = \omega_{21}$) capita che $|C_1(t)$ può essere 1, ovvero i livelli energetici vengono periodicamente "svuotati". In questo modo si hanno fenomeni interessanti quali il funzionamento dei laser.